% ID: SEL-D0048-OB-E03
% ESTADO: cotejada. Segunda pasada visual y borrador compilado revisados.
% Fuente: Selectividad/sel_2013_matematicas_aplicadas/junio/reserva a/EXAMEN 5 -ANDALUCÍA.pdf
% SHA-256: 76b6f6a138ec0917a2e3b3a2925bccca739af3363e419603b6814d835860441e
% Página: 2; opción B; ejercicio original 3.

A la Junta General de Accionistas de una empresa asisten 105 accionistas de los cuales 45 tienen menos de 40 años y 18 más de 60 años. Sometida a votación una propuesta, es rechazada por la tercera parte de los menores de 40 años, por la tercera parte de los que están entre 40 y 60 años y por 4 personas mayores de 60 años; los demás la aceptan.
\begin{enumerate}[label=\alph*)]
\item \textbf{(0.75 puntos)} Calcule la probabilidad de que, elegida una persona al azar, tenga menos de 40 años y haya aceptado la propuesta.
\item \textbf{(0.75 puntos)} La prensa afirmó que la propuesta había sido aceptada por el 80\% de los asistentes, \textquestiondown{}es correcta la afirmación?
\item \textbf{(1 punto)} Si una persona escogida al azar ha rechazado la propuesta, \textquestiondown{}qué probabilidad hay de que tenga más de 60 años?
\end{enumerate}
